% TOP_PROBLEM: {"id": 2828, "title": "Kirby Problem 3.30", "category_id": 7, "category": {"id": 7, "name": "topology", "display_name": "Topology", "description": "Properties preserved under continuous deformations.", "slug": "topology", "order_index": 7, "created_at": "2026-07-31T15:26:25.671Z"}, "status": "open", "problem_number": "KP-3.30", "set_id": 11}
% FIRST_POSTED: 2026-10-01
% ENTRY_KIND: statement_only
% UnsolvedMath ID 2828; source catalogue code KP-3.30
% !TeX program = lualatex
% Definition and sources only; this file is not an LLM solution attempt.
\documentclass[11pt,a4paper]{article}
\usepackage[margin=25mm]{geometry}
\usepackage{fontspec}
\setmainfont{lmroman10-regular.otf}[BoldFont=lmroman10-bold.otf,
 ItalicFont=lmroman10-italic.otf,BoldItalicFont=lmroman10-bolditalic.otf]
\usepackage{amsmath,amssymb,mathtools}
\usepackage{unicode-math}
\setmathfont{latinmodern-math.otf}
\AtBeginDocument{\let\setminus\smallsetminus}
\usepackage{xurl,hyperref}
\hypersetup{colorlinks=true,urlcolor=blue}
\setcounter{secnumdepth}{0}
\setlength{\parindent}{0pt}
\setlength{\parskip}{5pt}
\setlength{\emergencystretch}{3em}
\begin{document}
\section*{Kirby Problem 3.30}\label{2828}
{\small UnsolvedMath ID 2828 · KP-3.30 · Topology · Kirby's Problems in Low-Dimensional Topology}
\subsection{Definitions and mathematical statement}
Let $M$ be a compact orientable three-manifold specified by a triangulation
with $t$ tetrahedra. In the finite-volume convention, $M$ is hyperbolic if
its interior admits a complete sectional-curvature $-1$ metric of finite
volume; the closed case is included, and nonempty cusp boundary consists
of tori. This convention is different from prescribing geodesic boundary.

(a) Determine the bit complexity of deciding whether $M$ is hyperbolic.
(b) For a positive instance, determine the complexity of constructing a
finite, verifiable description of its hyperbolic structure.
(c) Can the decision problem lie in NP and the construction problem in FNP?

For precision, membership in NP requires polynomial-length certificates
whose correctness a deterministic algorithm checks in time polynomial in
the triangulation input length. The FNP version asks for a polynomially
balanced, polynomial-time checkable relation between the input and a finite
description of a structure. Such a description must specify how exact
geometric parameters are encoded and certified, for example algebraic
parameters with isolating data and verified gluing and completeness
conditions in a suitable geometric decomposition.

Uncertified floating-point approximations are not such certificates. A
solution should state its output representation and prove both geometric
validity and the claimed size and verification-time bounds.
\subsection{Short English statement}
How difficult is it to recognize and construct a hyperbolic structure from a three-manifold triangulation, and can valid structures have efficiently checkable certificates?
\subsection{Sources}
\begin{itemize}
\item[S1] ULAM.AI and the UnsolvedMath Contributors, supplied problems.json, record 2828 (KP-3.30), Kirby's Problems in Low-Dimensional Topology. Catalogue identity and source transcription; CC BY 4.0.
\url{https://huggingface.co/datasets/ulamai/UnsolvedMath}.
\item[S2] K3: A New Problem List in Low-Dimensional Topology, Problem 3.30. Expanded formulation of the supplied catalogue statement.
\url{https://math.berkeley.edu/sites/default/files/surv-295-ruberman-watermarked-author-pdf.pdf}.
\end{itemize}
\end{document}
